\documentclass[oneside,numbers,spanish]{ezthesis}
%% # Opciones disponibles para el documento #
%%
%% Las opciones con un (*) son las opciones predeterminadas.
%%
%% Modo de compilar:
%%   draft            - borrador con marcas de fecha y sin im'agenes
%%   draftmarks       - borrador con marcas de fecha y con im'agenes
%%   final (*)        - version final de la tesis
%%
%% Tama'no de papel:
%%   letterpaper (*)  - tama'no carta (Am'erica)
%%   a4paper          - tama'no A4    (Europa)
%%
%% Formato de impresi'on:
%%   oneside          - hojas impresas por un solo lado
%%   twoside (*)      - hijas impresas por ambos lados
%%
%% Tama'no de letra:
%%   10pt, 11pt, o 12pt (*)
%%
%% Espaciado entre renglones:
%%   singlespace      - espacio sencillo
%%   onehalfspace (*) - espacio de 1.5
%%   doublespace      - a doble espacio
%%
%% Formato de las referencias bibliogr'aficas:
%%   numbers          - numeradas, p.e. [1]
%%   authoryear (*)   - por autor y a'no, p.e. (Newton, 1997)
%%
%% Opciones adicionales:
%%   spanish         - tesis escrita en espa'nol
%%
%% Desactivar opciones especiales:
%%   nobibtoc   - no incluir la bibiolgraf'ia en el 'Indice general
%%   nofancyhdr - no incluir "fancyhdr" para producir los encabezados
%%   nocolors   - no incluir "xcolor" para producir ligas con colores
%%   nographicx - no incluir "graphicx" para insertar gr'aficos
%%   nonatbib   - no incluir "natbib" para administrar la bibliograf'ia

%% Paquetes adicionales requeridos se pueden agregar tambi'en aqu'i.
%% Por ejemplo:
%\usepackage{subfig}
%\usepackage{multirow}

%% # Datos del documento #
%% Nota que los acentos se deben escribir: \'a, \'e, \'i, etc.
%% La letra n con tilde es: \~n.

\author{Jes\'us Emil Ugueto P\'erez y Marvin Gonzalez}
\title{Desarrollo de un Juego Serio para apoyo a la enseñanza de la lectura mediante el m\'etodo mixto}
\degree{Licenciatura en Ciencias de la Computaci\'on}
\supervisor{Nombre de mi Asesor}
\institution{Universidad Central de Venezuela}
\faculty{Facultad de Ciencias}
\department{Escuela de Computaci\'on}

%% # M'argenes del documento #
%% 
%% Quitar el comentario en la siguiente linea para austar los m'argenes del
%% documento. Leer la documentaci'on de "geometry" para m'as informaci'on.

%\geometry{top=40mm,bottom=33mm,inner=40mm,outer=25mm}

%% El siguiente comando agrega ligas activas en el documento para las
%% referencias cruzadas y citas bibliogr'aficas. Tiene que ser *la 'ultima*
%% instrucci'on antes de \begin{document}.
\hyperlinking
\begin{document}

%% En esta secci'on se describe la estructura del documento de la tesis.
%% Consulta los reglamentos de tu universidad para determinar el orden
%% y la cantidad de secciones que debes de incluir.

%% # Portada de la tesis #
%% Mirar el archivo "titlepage.tex" para los detalles.
%% ## Construye tu propia portada ##
%% 
%% Una portada se conforma por una secuencia de "Blocks" que incluyen
%% piezas individuales de informaci'on. Un "Block" puede incluir, por
%% ejemplo, el t'itulo del documento, una im'agen (logotipo de la universidad),
%% el nombre del autor, nombre del supervisor, u cualquier otra pieza de
%% informaci'on.
%%
%% Cada "Block" aparece centrado horizontalmente en la p'agina y,
%% verticalmente, todos los "Blocks" se distruyen de manera uniforme 
%% a lo largo de p'agina.
%%
%% Nota tambi'en que, dentro de un mismo "Block" se pueden cortar
%% lineas usando el comando \\
%%
%% El tama'no del texto dentro de un "Block" se puede modificar usando uno de
%% los comandos:
%%   \small      \LARGE
%%   \large      \huge
%%   \Large      \Huge
%%
%% Y el tipo de letra se puede modificar usando:
%%   \bfseries - negritas
%%   \itshape  - it'alicas
%%   \scshape  - small caps
%%   \slshape  - slanted
%%   \sffamily - sans serif
%%
%% Para producir plantillas generales, la informaci'on que ha sido inclu'ida
%% en el archivo principal "tesis.tex" se puede accesar aqu'i usando:
%%   \insertauthor
%%   \inserttitle
%%   \insertsupervisor
%%   \insertinstitution
%%   \insertdegree
%%   \insertfaculty
%%   \insertdepartment
%%   \insertsubmitdate

\begin{titlepage}
  \TitleBlock{\includegraphics[height=4cm]{img/logoucv}}
  \TitleBlock{\scshape\insertinstitution}
  \TitleBlock[\bigskip]{\scshape\insertfaculty}
  \TitleBlock{\Huge\scshape\inserttitle}
  \TitleBlock{\scshape
    Tesis presentada por \insertauthor \\
    para obtener el grado de \insertdegree}
  \TitleBlock{\insertsubmitdate}
  \TitleBlock[\bigskip]{\insertdepartment}
\end{titlepage}

%% Nota 1:
%% Se puede agregar un escudo o logotipo en un "Block" como:
%%   \TitleBlock{\includegraphics[height=4cm]{escudo_uni}}
%% y teniendo un archivo "escudo_uni.pdf", "escudo_uni.png" o "escudo_uni.jpg"
%% en alg'un lugar donde LaTeX lo pueda encontrar.

%% Nota 2:
%% Normalmente, el espacio entre "Blocks" se extiende de modo que el
%% contenido se reparte uniformemente sobre toda la p'agina. Este
%% comportamiento se puede modificar para mantener fijo, por ejemplo, el
%% espacio entre un par de "Blocks". Escribiendo:
%%   \TitleBlock{Bloque 1}
%%   \TitleBlock[\bigskip]{Bloque2}
%% se deja un espacio "grande" y de tama~no fijo entre el bloque 1 y 2.
%% Adem'as de \bigskip est'an tambi'en \smallskip y \medskip. Si necesitas
%% aun m'as control puedes usar tambi'en, por ejemplo, \vspace*{2cm}.




%% # Prefacios #
%% Por cada prefacio (p.e. agradecimientos, resumen, etc.) crear
%% un nuevo archivo e incluirlo aqu'i.
%% Para m'as detalles y un ejemplo mirar el archivo "gracias.tex".
%% Las secciones del "prefacio" inician con el comando \prefacesection{T'itulo}
%% Este tipo de secciones *no* van numeradas, pero s'i aparecen en el 'indice.
%%
%% Si quieres agregar una secci'on que no vaya n'umerada y que *tampoco*
%% aparesca en el 'indice, usa entonces el comando \chapter*{T'itulo}
%%
%% Recuerda que aqu'i ya puedes escribir acentos como: 'a, 'e, 'i, etc.
%% La letra n con tilde es: 'n.

\prefacesection{Agradecimientos}


%% # 'Indices y listas de contenido #
%% Quitar los comentarios en las lineas siguientes para obtener listas de
%% figuras y cuadros/tablas.
\tableofcontents
%\listoffigures
%\listoftables

%% # Cap'itulos #
%% Por cada cap'itulo hay que crear un nuevo archivo e incluirlo aqu'i.
%% Mirar el archivo "intro.tex" para un ejemplo y recomendaciones para
%% escribir.
%% Los cap'itulos inician con \chapter{T'itulo}, estos aparecen numerados y
%% se incluyen en el 'indice general.
%%
%% Recuerda que aqu'i ya puedes escribir acentos como: 'a, 'e, 'i, etc.
%% La letra n con tilde es: 'n.

%Existen dos tipos de citas bibliograf'icas: usa \verb|\citep{..}| para
%citas en \emph{par'entesis} y \verb|\citet{..}| para citas
%en el \emph{texto}. Por ejemplo, estudios reciente han mostrado nuevos e
%interesantes modelos que se pueden aplicar para reformular teor'ias
%f'isicas~\citep{NewCam97}. Mientras que, el trabajo de \citet{Rofl06} fue
%considerado muy divertido por una significativa fracci'on de la comunidad
%de investigadores. Tambi'en es posible citar a varios trabajos en una sola
%referencia \citep{Lamport86,Knuth84}.
%
%Estos comandos para producir citas bibliograficas son provistos por
%el paquete \textsf{natbib}. Para obtener m'as informaci'on, consulta la
%documentaci'on de ese paquete~\citep{doc:natbib}. Por su parte, en
%la documentaci\'on de \textsf{geometry} puedes encontrar detalles
%adicionales sobre el sistema para ajustar los m'argenes del
%documento~\citep{doc:geometry}. Lo que sigue
%es un mont'on de texto sin sentido en lat'in que utilizaremos para llenar
%algunas p'aginas.



\prefacesection{Introducci'on}


%Learning to read is an important educational goal. For both
%children and adults, the ability to read opens up new worlds
%and opportunities. It enables us to gain new knowledge, enjoy
%literature, and do everyday things that are part and parcel of
%modern life

%a nivel universal, se consideran tres
%aprendizajes esenciales para la vida: la lectura, la escritura y el pensamiento lógico-matemático.

%% Reading development involves the participation of children, parents,
%% educators and the community as a whole. 

%% A child’s initial contact with words and symbols happens before going to
%% school.


\chapter{MARCO TE'ORICO}

%% --------------------- RESUMEN DEL CAPITULO ---------------------


\section{Alfebitzaci'on}

%%ALFABETIZACION Y EDUCACION - UNESCO \infante13

% --Alfabetizacion funcional:  “está alfabetizada toda persona que puede leer y escribir, comprendiéndola,
% una breve y sencilla exposición de hechos relativos a su vida cotidiana”.
% “El objetivo de la alfabetización funcional estaba determinado por la urgencia de
% movilizar, formar y educar la mano de obra aún subutilizada, para volverla más
% productiva, más útil a ella misma y a la sociedad” 

% Las ideas progresistas se incorporan al Simposio Internacional de Alfabetización
% celebrado en Persépolis 1975. En la declaración final, se define la “alfabetización
% no solo como el aprendizaje de la lectura, la escritura y el cálculo, sino como una
% contribución a la liberación de la persona y a su pleno desarrollo. Así concebida, la
% alfabetización crea las condiciones para la adquisición de una conciencia crítica de
% las contradicciones y los objetivos de la sociedad en que se vive; también estimula
% su iniciativa y participación en la creación de un proyecto capaz de actuar en el
% mundo, de transformarlo, y de definir los objetivos de un auténtico desarrollo
% humano”(IIALM, 1977, p. 636).

% -- La alfabetización entendida como un continuum de habilidades
%  “la habilidad de entender y utilizar información impresa en actividades cotidianas en el hogar,
% la comunidad y el trabajo”.

% -Hacia el concepto de alfabetización como puerta de entrada al aprendizaje permanente
% La alfabetización comienza a ser vista como un cimiento fundamental de todos los aprendizajes.

El concepto de alfabetizaci'on ha ido evolucionando a trav'es del tiempo. La estructura de la
sociedad est'a asociada historicacmente con la alfabetizaci'on, por tanto, los cambios culturales,
pol'iticos, y econ'omicos que sufre la sociedad afectan directamente la definici'on de alfabeticaci'on.

Hacia los a'nos 50 se ten'ia el concepto tradicional de alfabetizaci'on, y se consideraba alfabeta 
a aquellas personas que supieran decodificar los signos necesarios para leer y escribir \cite{infante13}, sin tomar en cuenta el uso que se le pudiera dar a estas habilidades.

\subsection{Alfabetizaci'on funcional}
En la d'ecada de los 60 se introduce el concepto de ``\textbf{alfabetizaci'on funcional}" que se definida por \citet{Unesco1970} como ``cualquier operaci'on de alfabetizaci'on concebida como un componente de los proyectos de desarrollo econom'omico y social". Se separa de la alfabetizaci'on tradicional al enlazar el aprendizaje de la lectura y la escritura con la formaci'on de los individuos como profesionales, y entiende que estas habilidades permiten la adquisici'on de aptitudes profesionales y conocimientos empleables en 'ambitos concretos. Sin embargo este enfoque inicial estaba orientado a formar y educar la mano de obra con la intenci'on de hacerla m'as productiva. M'as adelante las ideas progresistas comienzan a vincular la alfabetizaci'on con la concientizaci'on y el cambio social. A finales de la d'ecada de los 70 se incorpora una nueva definici'on de alfabetizaci'on funcional: ``Es alfabeto funcional la persona que puede emprender aquellas actividades en que la alfabetizaci'on es necesaria para la actuaci'on eficaz en su grupo y comunidad y que le permitan asimismo seguir valie'endose de la lectura, la escritura y la aritm'etica al servicio de su propio desarrollo y del desarrollo de la comunidad". {\color{red} Citar Acta UNESCO 1978}

\subsection{Alfabetizaci'on como puerta de entrada al aprendizaje permanente}



\section{La educaci'on inicial}

%% Venezuela AEPI
% La atención educativa e integral a la primera infancia, se define como Educación
% Inicial, la cual se concibe como el proceso de desarrollo y aprendizaje de los niños
% niñas, como un continuo que se inicia desde la gestación hasta los 6 años de edad o
% cuando se incorpore a la educación primaria y tiene como finalidad propiciar el
% máximo desarrollo de las potencialidades socio-afectivas, cognitivas, lingüísticas,
% motoras y físicas; considerando sus experiencias socio-educativas, intereses y
% necesidades.

\section{Influencia de las TIC en la educaci'on inicial}


\section{Juegos serios}






\chapter{TRABAJOS PREVIOS}

\chapter{PROPUESTA DE TESIS}
%%The prevalence of electronic media in the lives of young
%%children means that they are spending an increasing number
%%of hours per week in front of and engaged with screens of all
%%kinds, including televisions, computers, smartphones, tablets,
%%handheld game devices, and game consoles

%%With guidance, these various
%%technology tools can be harnessed for learning and development;
%%without guidance, usage can be inappropriate and/or
%%interfere with learning and development.

%% A DAY IN THE LIFE OF A ECE TEACHER
% The issues were identified as “insufficient time”, a “lack of resources” and also “working with
% parents” and these are discussed below

\section{Planteamiento del problema}

%% Alfabetizacion inicial
% un aprendizaje lector insuficiente durante los primeros años de la educación básica 
% limita los niveles de procesamiento cognitivo y el acceso al conocimiento, 
% en las diversas áreas de contenido de los programas escolares.

% recientemente se ha aceptado que muchos de los procesos más importantes para el 
% desarrollo de la alfabetización tienen lugar antes de que el alumno entre por
% primera vez en un aula (Pressley, 1999). 

El estado de desarrollo de las sociedades y los individuos que la componen est'an fuertemente
asociado. Una sociedad con un alto nivel de pobreza es indicativa de un bajo acceso a la educaci'on, 
y se vincula con la reproducci'on del analfabetismo. Seg'un \citet{infante13} "no hay posibilidad
e alcanzar una democracia efectiva, mientras gran parte de la poblaci'on se mantenga fuera del
acceso a la lengua escrita". La alfabetizaci'on es entonces un elemento fundamental para el debido funcionamiento de la sociedad, ha de ser vista como un derecho universal, y es deber de la sociedad misma promoverla. Sin embargo, se debe entender la alfabetizaci'on no como un proceso
simple que s'olo comprende leer y escribir, sino como a la puerta de entrada al aprendizaje
permanente. 

{\color{red} Necesito un p'arrafo de puente. Debe hablar sobre la educaci'on inicial,
el aprendizaje de la lectura (metodos mixto), y como pueden las TICs ayudar en esto.}

El r'apido avance de las Tecnologías de la Información y la Comunicación (TIC) influye
significativamente sobre diferentes 'ambitos sociales, como en la econom'ia, la cultura, y 
la educaci'on. De acuerdo con \cite{CurriculoSEPB} los egresados y las egresadas del subsistema
de Educaci'on Primaria Bolivariana debe poseer entre sus caracter'iasticas el uso
de las TIC en el proceso de aprendizaje, as'i mismo 'estas forman parte de los ejes
integradores del subsistema, lo que significa que todos los procesos educativos del subsistema
deben considerar las TIC para orientar las experiencias del aprendizaje, as'i como fomentar valores,
actitudes y virtudes.

Para la incorporaci'on de las TIC en la educaci'on se debe tomar en cuenta, entre otras cosas,
la cantidad, calidad y costos de software desarrollado espec'ificamente con fines educativos, y
enfoncados a ni'nos y ni'nas de corta edad, o c'omo el uso de las computadoras ha sido asociado 
comunmente con el entretenimiento y no con actividades que desarrollen las destrezas b'asicas 
necesarias para un desenvolvimiento 'optimo \cite{Archila}.

Favorecer el uso de las TIC en la educaci'on depende entonces de la accesibilidad que se le de
a los miembros de la comunidad a las herramientas de hardware (computadoras, tablets, tel'efonos
m'oviles), y software (programas, aplicaciones, videojuegos). Lo primero es atacado por iniciativas
como el Proyecto Canaima Educativo {\color{red} (incluir referencia)}. Lo segundo implica el desarrollo de herramientas bien fundamentadas en las diferentes metodolog'ias del aprendizaje,
 basadas en licencias libres, y con un alto nivel de portabilidad, mantenibilidad y extensibilidad.
 
El cambio curricular que incluye a las TIC como eje integrador del subsistema de Educaci'on
Primaria Bolivariana data del a~no 2007, por lo que gran parte del software que se utiliza
es el que ya exist'ia para entonces. Esto significa que no hay software actualizado, acorde con 
las nuevas tecnolog'ias, algo fundamental un entorno que evoluciona tan velozmente como lo son
las TIC. As'i mismo hay que tomar en cuenta que en la actualidad los ni'nos y ni'nas nacen, y se
desarrollan de la mano de la tecnolog'ia, est'an fuertemente adaptados a esta, y siguen 
f'acilmente los patrones y tendencias actuales, llegando incluso a rechazar aquello que consideren
desfasado.

{\color{red} Hace falta un p'arrafo de cierre.}

\section{Objetivo General}

Desarrollar una herramienta que permita {\color{red}emplear} el m'etodo mixto de ense'nanza de la lectura, enfocada a ni'nos y ni'nas durante la etapa de educaci'on inicial, que siga las tendencias
tecnol'ogicas actuales, y sea altamente portable, mantenible y extensible. 


\section{Objetivos Espec'ificos}


\section{Justificaci'on de la propuesta}




%\include{conclu}

\appendix
%% Cap'itulos incluidos despues del comando \appendix aparecen como ap'endices
%% de la tesis.
%\include{apendiceA}
%\include{apendiceB}
%\include{apendiceC}

%% Incluir la bibliograf'ia. Mirar el archivo "biblio.bib" para m'as detales
%% y un ejemplo.
\bibliography{biblio}

\end{document}
